\ifdefined\gkver\else\def\gkver{student}\fi
\documentclass[\gkver]{gaokaozhenti}
\begin{document}

\chapter{2015年海南理综}
%% paperType: 理综物理部分

\begin{enumerate}
\item
%% number: 1
%% paperName: 2015年海南理综第 1 题 3 分
%% typeId: 1
%% type: 单选题
%% chapter: 磁场
%% point: 洛伦兹力
%% method: 无
%% score: 3
%% degree: 650
%% duplicateId: 0
%% body:
如图，a 是竖直平面 P 上的一点。P 前有一条形磁铁垂直于 P，且 S 极朝向 a 点，P 后一电子在偏转线圈和条形磁铁的磁场的共同作用下，在水平面内向右弯曲经过 a 点。在电子经过 a 点的瞬间，条形磁铁的磁场对该电子的作用力的方向\xzanswer{A}

\onepicture[h=3.0cm]{figs/01.png}

\fourchoices[answer=A]
{向上}
{向下}
{向左}
{向右}

%% answer: A
%% memo:
\memoanswer{
	条形磁铁的磁感线方向在 a 点为垂直 P 向外，粒子在条形磁铁的磁场中向右运动，所以根据左手定则可得电子受到的洛伦兹力方向向上，选项 A 正确。
}

\item
%% number: 2
%% paperName: 2015年海南理综第 2 题 3 分
%% typeId: 1
%% type: 单选题
%% chapter: 电磁感应
%% point: 切割磁感线电动势
%% method: 等效替代
%% score: 3
%% degree: 650
%% duplicateId: 0
%% body:
如图，空间有一匀强磁场，一直金属棒与磁感应强度方向垂直，当它以速度 $v$ 沿与棒和磁感应强度都垂直的方向运动时，棒两端的感应电动势大小为 $\varepsilon$；将此棒弯成两段长度相等且相互垂直的折弯，置于磁感应强度相垂直的平面内，当它沿两段折线夹角平分线的方向以速度 $v$ 运动时，棒两端的感应电动势大小为 $\varepsilon'$，则 $\frac{\varepsilon'}{\varepsilon}$ 等于\xzanswer{B}

\onepicture[h=3.0cm]{figs/02.png}

\fourchoices[answer=B]
{$\frac{1}{2}$}
{$\frac{\sqrt{2}}{2}$}
{1}
{$\sqrt{2}$}

%% answer: B
%% memo:
\memoanswer{
	设折弯前导体切割磁感线的长度为 $l$，产生的感应电动势为 $\varepsilon=Blv$；折弯后，导体切割磁场的有效长度为 $\frac{\sqrt{2}}{2}l$，产生的感应电动势为 $\varepsilon'=B\cdot\frac{\sqrt{2}}{2}lv=\frac{\sqrt{2}}{2}\varepsilon$，所以 $\frac{\varepsilon'}{\varepsilon}=\frac{\sqrt{2}}{2}$，选项 B 正确。
}

\item
%% number: 3
%% paperName: 2015年海南理综第 3 题 3 分
%% typeId: 1
%% type: 单选题
%% chapter: 功和能
%% point: 交通工具启动
%% method: 无
%% score: 3
%% degree: 650
%% duplicateId: 0
%% body:
假设摩托艇受到的阻力的大小正比于它的速率。如果摩托艇发动机的输出功率变为原来的 2 倍，则摩托艇的最大速率变为原来的\xzanswer{D}

\fourchoices[answer=D]
{4 倍}
{2 倍}
{$\sqrt{3}$ 倍}
{$\sqrt{2}$ 倍}

%% answer: D
%% memo:
\memoanswer{
	设 $F_f=kv$，当阻力等于牵引力时，速度最大，输出功率变化前，有 $P=Fv=F_fv=kv^2$，变化后有 $2P=F'v'=F_f'v'=kv'^2$，联立解得 $v'=\sqrt{2}v$，选项 D 正确。
}

\item
%% number: 4
%% paperName: 2015年海南理综第 4 题 3 分
%% typeId: 1
%% type: 单选题
%% chapter: 曲线运动
%% point: 竖直平面圆周运动
%% method: 无
%% score: 3
%% degree: 450
%% duplicateId: 0
%% body:
如图，一半径为 $R$ 的半圆形轨道竖直固定放置，轨道两端等高；质量为 $m$ 的质点自轨道端点 P 由静止开始滑下，滑到最低点 Q 时，对轨道的正压力为 $2mg$，重力加速度大小为 $g$。质点自 P 滑到 Q 的过程中，克服摩擦力所做的功为\xzanswer{C}

\onepicture[h=3.2cm]{figs/04.png}

\fourchoices[answer=C]
{$\frac{1}{4}mgR$}
{$\frac{1}{3}mgR$}
{$\frac{1}{2}mgR$}
{$\frac{\pi}{4}mgR$}

%% answer: C
%% memo:
\memoanswer{
	在 Q 点质点受到重力和支持力作用，其合力提供向心力，有 $F_N-mg=m\frac{v^2}{R}$，$F_N=2mg$，联立解得 $v=\sqrt{Rg}$，下落过程中重力做正功，摩擦力做负功，根据动能定理可得 $mgR-W_f=\frac{1}{2}mv^2$，解得 $W_f=\frac{1}{2}mgR$，所以克服摩擦力做功 $\frac{1}{2}mgR$，选项 C 正确。
}

\item
%% number: 5
%% paperName: 2015年海南理综第 5 题 3 分
%% typeId: 1
%% type: 单选题
%% chapter: 电场
%% point: 带电粒子的直线运动
%% method: 无
%% score: 3
%% degree: 650
%% duplicateId: 0
%% body:
如图，一充电后的平行板电容器的两极板相距 $l$。在正极板附近有一质量为 $M$、电荷量为 $q$（$q>0$）的粒子；在负极板附近有另一质量为 $m$、电荷量为 $-q$ 的粒子。在电场力的作用下，两粒子同时从静止开始运动。已知两粒子同时经过一平行于正极板且与其相距 $\frac{2}{5}l$ 的平面。若两粒子间相互作用力可忽略，不计重力，则 $M:m$ 为\xzanswer{A}

\onepicture[h=3.0cm]{figs/05.png}

\fourchoices[answer=A]
{$3:2$}
{$2:1$}
{$5:2$}
{$3:1$}

%% answer: A
%% memo:
\memoanswer{
	设电场强度为 $E$，两粒子的运动时间相同，对 $M$ 有：$\frac{2}{5}l=\frac{1}{2}\frac{qE}{M}t^2$；对 $m$ 有：$\frac{3}{5}l=\frac{1}{2}\frac{qE}{m}t^2$，联立解得 $\frac{M}{m}=\frac{3}{2}$，选项 A 正确。
}

\item
%% number: 6
%% paperName: 2015年海南理综第 6 题 3 分
%% typeId: 1
%% type: 单选题
%% chapter: 曲线运动
%% point: 平抛运动
%% method: 无
%% score: 3
%% degree: 450
%% duplicateId: 0
%% body:
若在某行星和地球上相对于各自水平地面附近相同的高度处、以相同的速率平抛一物体，它们在水平方向运动的距离之比为 $2:\sqrt{7}$。已知该行星质量约为地球的 7 倍，地球的半径为 $R$。由此可知，该行星的半径为\xzanswer{C}

\fourchoices[answer=C]
{$\frac{1}{2}R$}
{$\frac{7}{2}R$}
{$2R$}
{$\frac{\sqrt{7}}{2}R$}

%% answer: C
%% memo:
\memoanswer{
	由平抛运动规律：$x=v_0t$，$h=\frac{1}{2}gt^2$ 得 $x=v_0\sqrt{\frac{2h}{g}}$，两种情况下，抛出的速度相同，高度相同，故 $\frac{g_{\text{行}}}{g_{\text{地}}}=\frac{7}{4}$；由 $G\frac{Mm}{R^2}=mg$ 可得 $g=\frac{GM}{R^2}$，故 $\frac{g_{\text{行}}}{g_{\text{地}}}=\frac{M_{\text{行}}/R_{\text{行}}^2}{M_{\text{地}}/R_{\text{地}}^2}=\frac{7}{4}$，解得 $R_{\text{行}}=2R$，选项 C 正确。
}

\item
%% number: 7
%% paperName: 2015年海南理综第 7 题 5 分
%% typeId: 2
%% type: 多选题
%% chapter: 电场
%% point: 电势能电势差电场力做功
%% method: 无
%% score: 5
%% degree: 450
%% duplicateId: 0
%% body:
如图，两电荷量分别为 $Q$（$Q>0$）和 $-Q$ 的点电荷对称地放置在 $x$ 轴上原点 $O$ 的两侧，a 点位于 $x$ 轴上 $O$ 点与点电荷 $Q$ 之间，b 位于 $y$ 轴 $O$ 点上方，取无穷远处的电势为零。下列说法正确的是\xzanswer{BC}

\onepicture[h=3.8cm]{figs/07.png}

\fourchoices[answer=BC]
{b 点的电势为零，电场强度也为零}
{正的试探电荷在 a 点的电势能大于零，所受电场力方向向右}
{将正的试探电荷从 $O$ 点移到 a 点，必须克服电场力做功}
{将同一正的试探电荷先后从 $O$、b 点移到 a 点，后者电势能的变化较大}

%% answer: BC
%% memo:
\memoanswer{
	因为等量异种电荷在其连线的中垂线上的电场方向为水平指向负电荷，所以电场方向与中垂线方向垂直，故中垂线为等势线，因为中垂线延伸到无穷远处，所以中垂线的电势为零，故 b 点的电势为零，但是电场强度不为零，A 错误；等量异种电荷连线上，电场方向由正电荷指向负电荷，方向水平向右，在中点 $O$ 处电势为零，$O$ 点左侧电势为正，右侧电势为负，又知道正电荷在正电势处电势能为正，故 B 正确；$O$ 点的电势低于 a 点的电势，电场力做负功，所以必须克服电场力做功，C 正确；$O$ 点和 b 点的电势相等，所以先后从 $O$、b 点移到 a 点，电场力做功相等，电势能变化相同，D 错误。
}

\item
%% number: 8
%% paperName: 2015年海南理综第 8 题 5 分
%% typeId: 2
%% type: 多选题
%% chapter: 运动和力
%% point: 瞬时加速度
%% method: 无
%% score: 5
%% degree: 450
%% duplicateId: 0
%% body:
如图，物块 a、b 和 c 的质量相同，a 和 b、b 和 c 之间用完全相同的轻弹簧 $S_1$ 和 $S_2$ 相连，通过系在 a 上的细线悬挂于固定点 $O$。整个系统处于静止状态，现将细绳剪断，将物块 a 的加速度记为 $a_1$，$S_1$ 和 $S_2$ 相对原长的伸长分别为 $\Delta l_1$ 和 $\Delta l_2$，重力加速度大小为 $g$，在剪断瞬间\xzanswer{AC}

\onepicture[h=4.0cm]{figs/08.png}

\fourchoices[answer=AC]
{$a_1=3g$}
{$a_1=0$}
{$\Delta l_1=2\Delta l_2$}
{$\Delta l_1=\Delta l_2$}

%% answer: AC
%% memo:
\memoanswer{
	设物体的质量为 $m$，剪断细绳的瞬间，绳子的拉力消失，弹簧还没有来得及改变，所以剪断细绳的瞬间 a 受到重力和弹簧 $S_1$ 的拉力 $T_1$，剪断前对 b、c 和弹簧组成的整体分析可知 $T_1=2mg$，故 a 受到的合力 $F=mg+T_1=3mg$，故加速度 $a_1=\frac{F}{m}=3g$，A 正确、B 错误；设弹簧 $S_2$ 的拉力为 $T_2$，则 $T_2=mg$，根据胡克定律 $F=k\Delta x$ 可得 $\Delta l_1=2\Delta l_2$，C 正确、D 错误；选项 AC 正确。
}

\item
%% number: 9
%% paperName: 2015年海南理综第 9 题 5 分
%% typeId: 2
%% type: 多选题
%% chapter: 运动和力
%% point: 超重失重
%% method: 无
%% score: 5
%% degree: 450
%% duplicateId: 0
%% body:
如图，升降机内有一固定斜面，斜面上放一物体。开始时升降机做匀速运动，物块相对斜面匀速下滑，当升降机加速上升时，\xzanswer{BD}

\onepicture[h=3.2cm]{figs/09.png}

\fourchoices[answer=BD]
{物块与斜面间的摩擦力减小}
{物块与斜面间的正压力增大}
{物块相对于斜面减速下滑}
{物块相对于斜面匀速下滑}

%% answer: BD
%% memo:
\memoanswer{
	当升降机加速上升时，物体有竖直向上的加速度，则物块与斜面间的正压力增大，根据滑动摩擦力公式 $F_f=\mu F_N$ 可知接触面间的正压力增大，物体与斜面间的摩擦力增大，故 A 错误、B 正确；设斜面的倾角为 $\theta$，物体的质量为 $m$，当匀速运动时有 $mg\sin\theta=\mu mg\cos\theta$，即 $\sin\theta=\mu\cos\theta$，假设物体以加速度 $a$ 向上运动时，有 $F_N=m(g+a)\cos\theta$，$F_f=\mu m(g+a)\cos\theta$，因为 $\sin\theta=\mu\cos\theta$，所以 $m(g+a)\sin\theta=\mu m(g+a)\cos\theta$，故物体仍做匀速下滑运动，C 错误、D 正确；选项 BD 正确。
}

\item
%% number: 10
%% paperName: 2015年海南理综第 10 题 5 分
%% typeId: 2
%% type: 多选题
%% chapter: 交流电
%% point: 变压器
%% method: 无
%% score: 5
%% degree: 450
%% duplicateId: 0
%% body:
如图，一理想变压器原、副线圈匝数比为 $4:1$，原线圈与一可变电阻串联后，接入一正弦交流电源；副线圈电路中固定电阻的阻值为 $R_0$，负载电阻的阻值 $R=11R_0$，V 是理想电压表。现将负载电阻的阻值减小为 $R=5R_0$，保持变压器输入电流不变，此时电压表读数为 $5.0\UV$，则\xzanswer{AD}

\onepicture[h=3.0cm]{figs/10.png}

\fourchoices[answer=AD]
{此时原线圈两端电压的最大值约为 $34\UV$}
{此时原线圈两端电压的最大值约为 $24\UV$}
{原线圈两端原来的电压有效值约为 $68\UV$}
{原线圈两端原来的电压有效值约为 $48\UV$}

%% answer: AD
%% memo:
\memoanswer{
	当负载电阻的阻值减小为 $R=5R_0$ 时，由串并联电路规律，$R$ 两端电压为 $R_0$ 两端电压的 5 倍，因为电压表测量 $R$ 两端的电压，所以 $U_{R_0}=\frac{1}{5}\times5\UV=1\UV$，副线圈两端电压为 $U_2=6\UV$，根据公式 $\frac{U_1}{U_2}=\frac{n_1}{n_2}$ 可得此时原线圈两端电压的有效值 $U_1=24\UV$，所以此时原线圈两端电压的最大值约为 $24\sqrt{2}\UV\approx34\UV$，A 正确、B 错误；因为变化过程中变压器的输入电流不变，所以对于副线圈中变化前后电流也不变，则变化后电压 $U_2=IR_0+5IR_0=6IR_0$，变化前 $U_2'=IR_0+11IR_0=12IR_0$，所以 $U_2'=2U_2=12\UV$，根据公式 $\frac{U_1}{U_2}=\frac{n_1}{n_2}$ 可得原线圈两端原来的电压有效值约为 $48\UV$，C 错误、D 正确；选项 AD 正确。
}

\item
%% number: 11
%% paperName: 2015年海南理综第 11 题 6 分
%% typeId: 4
%% type: 实验题
%% chapter: 直线运动
%% point: 游标卡尺和螺旋测微仪
%% method: 无
%% score: 6
%% degree: 650
%% duplicateId: 0
%% body:
某同学利用游标卡尺和螺旋测微器分别测量一圆柱体工件的直径和高度，测量结果如图（a）和（b）所示。该工件的直径为\tkanswer[1.4cm]{1.220}$\Ucm$，高度为\tkanswer[1.4cm]{6.861}$\Umm$。

\twopicture[numA={（a）}, numB={（b）}, hA=4.0cm, hB=3.4cm]
{figs/11a.jpg}
{figs/11b.jpg}

%% answer:
\jdanswer{
	$1.220\Ucm$，$6.861\Umm$（$6.859$ 或 $6.861$，同样给分）
}
%% memo:
\memoanswer{
	由题图（a）读得工件的直径为：$12\Umm+4\times0.05\Umm=12.20\Umm=1.220\Ucm$。由题图（b）读得工件的高度为 $6.5\Umm+0.360\Umm=6.860\Umm$。
}

\item
%% number: 12
%% paperName: 2015年海南理综第 12 题 9 分
%% typeId: 4
%% type: 实验题
%% chapter: 电路
%% point: 电容
%% method: 无
%% score: 9
%% degree: 450
%% duplicateId: 0
%% body:
某同学利用图（a）所示电路测量电容器充电时两极板间的电压随时间的变化。实验中使用的器材为：电池 $E$（内阻很小）、开关 $S_1$ 和 $S_2$、电容器 $C$（约 $100\UuF$）、电阻 $R_1$（约 $200\UkO$）、电阻 $R_2$（$1\UkO$）、电压表 V（量程 $6\UV$）、秒表、导线若干。

\threepicture[numA={（a）}, numB={（b）}, numC={（c）}, hA=4.2cm, hB=4.2cm, hC=2.6cm]
{figs/12a.jpg}
{figs/12b.jpg}
{figs/12c.jpg}

\begin{enumerate}
	\item
	按图（a）所示的电路原理将图（b）中实物图连线。

	\item
	先闭合开关 $S_2$，再断开开关 $S_2$；闭合开关 $S_1$，同时按下秒表开始计时。若某时刻电压表示数如图（c）所示，电压表的读数为\tkanswer[1.4cm]{3.60}$\UV$（保留 2 位小数）。

	\item
	该同学每隔 $10\Us$ 记录一次电压表的读数 $U$，记录的数据如下表所示，在给出的坐标纸上绘出 $U-t$ 图线，已知只有一个数据点误差较大，该数据点对应的表中的时间是\tkanswer[1.2cm]{40}$\Us$。

\begin{center}
\begin{tabular}{|c|c|c|c|c|c|c|}
\hline
时间 $t/\Us$ & 10.0 & 20.0 & 30.0 & 40.0 & 50.0 & 60.0 \\
\hline
电压 $U/\UV$ & 2.14 & 3.45 & 4.23 & 4.51 & 5.00 & 5.18 \\
\hline
\end{tabular}
\end{center}

\onepicture[h=5.2cm]{figs/12d.png}

	\item
	电路中 $C$、$R_2$ 和 $S_1$ 构成的回路的作用是\tkanswer[4.0cm]{使充电的电容器很快放电，以便多次测量}。
\end{enumerate}

%% answer:
\jdanswer{
\begin{enumerate}
	\item
	连线如图所示。

	\item
	$3.60\UV$（$3.59$ 或 $3.61$，同样给分）

	\item
	$U-t$ 图线如图所示，$40\Us$

	\item
	使充电的电容器很快放电，以便多次测量
\end{enumerate}
}
%% memo:
\memoanswer{
	（2）电压表的量程为 $6\UV$，最小分度值为 $0.1\UV$，读数为 $3.60\UV$。

	（3）$U-t$ 图线如图所示，由图像可得一个数据点误差较大的是 $40\Us$。

	（4）电路中 $C$、$R_2$ 和 $S_1$ 构成的回路的作用使充电的电容器很快放电，以便多次测量。

	\onepicture[h=5.0cm]{figs/12e.jpg}

	\onepicture[h=5.0cm]{figs/12f.jpg}
}

\item
%% number: 13
%% paperName: 2015年海南理综第 13 题 10 分
%% typeId: 6
%% type: 计算题
%% chapter: 电磁感应
%% point: 电磁感应受力分析
%% method: 无
%% score: 10
%% degree: 650
%% duplicateId: 0
%% body:
如图，两平行金属导轨位于同一水平面上，相距 $l$，左端与一电阻 $R$ 相连；整个系统置于匀强磁场中，磁感应强度大小为 $B$，方向竖直向下。一质量为 $m$ 的导体棒置于导轨上，在水平外力作用下沿导轨以速度 $v$ 匀速向右滑动，滑动过程中始终保持与导轨垂直并接触良好。已知导体棒与导轨间的动摩擦因数为 $\mu$，重力加速度大小为 $g$。导轨和导体棒的电阻均可忽略。求：

\begin{enumerate}
	\item
	电阻 $R$ 消耗的功率；

	\item
	水平外力的大小。
\end{enumerate}

\onepicture[h=3.2cm, align=right]{figs/13.png}

%% answer:
\jdanswer{
\begin{enumerate}
	\item
	$\frac{B^2l^2v^2}{R}$

	\item
	$\mu mg+\frac{B^2l^2v}{R}$
\end{enumerate}
}
%% memo:
\memoanswer{
	（1）导体棒上的电动势为 $E=Bvl$①，

	电阻 $R$ 消耗的功率为 $P=\frac{E^2}{R}$②，

	联立①②式得 $P=\frac{(Bvl)^2}{R}$③。

	（2）导体棒所受到的导轨的摩擦力为 $f_1=\mu mg$④，

	所受到的磁场施加的安培力为 $f_2=IBl$⑤，

	式中 $I$ 为感应电流。由于导体棒做匀速直线运动，作用在导体棒上的外力 $F$ 应满足力的平衡条件

	$F=f_1+f_2$⑥，

	由欧姆定律得 $E=IR$⑦，

	联立⑤⑥⑦式得 $F=\mu mg+\frac{B^2l^2v}{R}$⑧。
}

\item
%% number: 14
%% paperName: 2015年海南理综第 14 题 13 分
%% typeId: 6
%% type: 计算题
%% chapter: 曲线运动
%% point: 平抛运动
%% method: 无
%% score: 13
%% degree: 450
%% duplicateId: 0
%% body:
如图，位于竖直水平面内的光滑轨道由四分之一圆弧 $ab$ 和抛物线 $bc$ 组成，圆弧半径 $Oa$ 水平，$b$ 点为抛物线顶点。已知 $h=2\Um$，$s=\sqrt{2}\Um$。取重力加速度大小 $g=10\Umsq$。

\begin{enumerate}
	\item
	一小环套在轨道上从 a 点由静止滑下，当其在 $bc$ 段轨道运动时，与轨道之间无相互作用力，求圆弧轨道的半径；

	\item
	若环从 $b$ 点由静止因微小扰动而开始滑下，求环到达 $c$ 点时速度的水平分量的大小。
\end{enumerate}

\onepicture[h=5.0cm, align=right]{figs/14.png}

%% answer:
\jdanswer{
\begin{enumerate}
	\item
	$0.25\Um$

	\item
	$\frac{2\sqrt{10}}{3}\Ums$
\end{enumerate}
}
%% memo:
\memoanswer{
	（1）一小环套在 $bc$ 段与轨道无相互作用，必然在该段以某一初速度 $v_b$ 做平抛运动，运动的轨迹与轨道相同。由平抛运动公式有

	$s=v_bt$①，

	$h=\frac{1}{2}gt^2$②，

	设圆弧轨道半径为 $R$，由机械能守恒定律得

	$mgR=\frac{1}{2}mv_b^2$③，

	联立①②③式，并代入题给条件得

	$R=0.25\Um$④。

	（2）环由 $b$ 处静止下滑过程中机械能守恒，设下滑至 $c$ 点的速度大小为 $v$，有

	$mgh=\frac{1}{2}mv^2$⑤，

	环在 $c$ 点速度的水平分量为

	$v_x=v\cos\theta$⑥，

	式中 $\theta$ 为环在 $c$ 点速度方向与水平方向的夹角。由题意知，环在 $c$ 点速度方向和以初速度 $v_b$ 做平抛运动的物体在 $c$ 点速度方向相同；而做平抛运动的物体末速度的水平分量为 $v_x'=v_b$，竖直分量 $v_y'$

	$v_y'=\sqrt{2gh}$⑦，

	因此 $\cos\theta=\frac{v_b}{\sqrt{v_b^2+v_y'^2}}$⑧，

	联立①②⑤⑥⑦⑧式得 $v_x=\frac{2\sqrt{10}}{3}\Ums$⑨。
}

\item
%% number: 15
%% paperName: 2015年海南理综第 15 题 8 分
%% typeId: 6
%% type: 计算题
%% chapter: 气体性质
%% point: 气体连接体
%% method: 无
%% score: 8
%% degree: 450
%% duplicateId: 0
%% body:
【物理——选修 3-3】（12 分）

（1）（4 分）已知地球大气层的厚度 $h$ 远小于地球半径 $R$，空气平均摩尔质量为 $M$，阿伏伽德罗常数为 $N_A$，地面大气压强为 $p_0$，重力加速度大小为 $g$。由此可以估算得，地球大气层空气分子总数为\tkanswer[2.8cm]{$\frac{4\pi N_Ap_0R^2}{Mg}$}，空气分子之间的平均距离为\tkanswer[2.8cm]{$\sqrt[3]{\frac{Mgh}{N_Ap_0}}$}。

（2）（8 分）如图，一底面积为 $S$、内壁光滑的圆柱形容器竖直放置在水平地面上，开口向上，内有两个质量均为 $m$ 的相同活塞 A 和 B；在 A 与 B 之间、B 与容器底面之间分别封有一定量的同样的理想气体，平衡时体积均为 $V$。已知容器内气体温度始终不变，重力加速度大小为 $g$，外界大气压强为 $p_0$。现假设活塞 B 发生缓慢漏气，致使 B 最终与容器底面接触，求活塞 A 移动的距离。

\onepicture[h=4.0cm, align=right]{figs/15.png}

%% answer:
\jdanswer{
\begin{enumerate}
	\item
	$\frac{4\pi N_Ap_0R^2}{Mg}$，$\sqrt[3]{\frac{Mgh}{N_Ap_0}}$（或 $\sqrt[3]{\frac{Mg[(R+h)^3-R^3]}{3R^2N_Ap_0}}$）

	\item
	$\frac{mgV}{(p_0S+mg)S}$
\end{enumerate}
}
%% memo:
\memoanswer{
	（1）设大气层中气体的质量为 $m$，由大气压强产生，$mg=p_0S$，而分子数 $n=\frac{mN_A}{M}$，地球表面积 $S=4\pi R^2$，故地球大气层空气分子总数 $n=\frac{4\pi N_Ap_0R^2}{Mg}$，假设每个分子占据一个小立方体，各小立方体紧密排列，则小立方体边长即为空气分子平均间距，设为 $a$，则 $a=\left(\frac{V}{n}\right)^{1/3}$；大气层中气体总体积为 $V=\frac{4}{3}\pi(R+h)^3-\frac{4}{3}\pi R^3$，所以 $a=\left(\frac{Mgh}{N_Ap_0}\right)^{1/3}$。

	（2）设平衡时，在 A 与 B 之间、B 与容器底面之间的气体压强分别为 $p_1$、$p_2$，由力的平衡条件有

	$p_1S=p_0S+mg$①，

	$p_2S=p_1S+mg$②，

	漏气发生后，设整个封闭气体体积为 $V'$，压强为 $p'$，由力的平衡条件有

	$p'S=p_0S+mg$③，

	由玻意耳定律得

	$p'V_1'=p_1V$④，$p'(V'-V_1')=p_2V$⑤，

	式中 $V_1'$ 是原来 A 与 B 之间的气体在漏气发生后所占的体积。

	设活塞 A 移动的距离为 $l$（取升高时为正），按几何关系有

	$V'=2V+Sl$⑥，

	联立①②③④⑤⑥式得

	$l=\left(\frac{mg}{p_0S+mg}\right)\frac{V}{S}$。
}

\item
%% number: 16
%% paperName: 2015年海南理综第 16 题 8 分
%% typeId: 6
%% type: 计算题
%% chapter: 光学
%% point: 光的折射
%% method: 无
%% score: 8
%% degree: 450
%% duplicateId: 0
%% body:
【物理——选修 3-4】（12 分）

（1）（4 分）一列沿 $x$ 轴正方向传播的简谐横波在 $t=0$ 时刻的波形如图所示，质点 P 的 $x$ 坐标为 $3\Um$。已知任意振动质点连续 2 次经过平衡位置的时间间隔为 $0.4\Us$。下列说法正确的是\xzanswer{BDE}（填正确答案标号。选对 1 个得 2 分，选对 2 个得 3 分，选 3 个得 4 分；每选错 1 个扣 2 分，最低得分为 0 分）

\onepicture[h=3.4cm]{figs/16a.png}

\fivechoices[answer=BDE]
{波速为 $4\Ums$}
{波的频率为 $1.25\UHz$}
{$x$ 坐标为 $15\Um$ 的质点在 $t=0.2\Us$ 时恰好位于波谷}
{$x$ 坐标为 $22\Um$ 的质点在 $t=0.2\Us$ 时恰好位于波峰}
{当质点 P 位于波峰时，$x$ 坐标为 $17\Um$ 的质点恰好位于波谷}

（2）（8 分）一半径为 $R$ 的半圆形玻璃砖，横截面如图所示。已知玻璃的全反射临界角 $\gamma$（$\gamma<\frac{\pi}{3}$）。与玻璃砖的底平面成（$\frac{\pi}{2}-\gamma$）角度、且与玻璃砖横截面平行的平行光射到玻璃砖的半圆柱面上。经柱面折射后，有部分光（包括与柱面相切的入射光）能直接从玻璃砖底面射出。若忽略经半圆柱内表面反射后射出的光，求底面透光部分的宽度。

\onepicture[h=4.0cm]{figs/16b.png}

%% answer:
\jdanswer{
\begin{enumerate}
	\item
	BDE

	\item
	$\frac{R}{2\cos\gamma}$
\end{enumerate}
}
%% memo:
\memoanswer{
	（1）任意振动质点连续 2 次经过平衡位置的时间间隔为 $0.4\Us$，则 $T=0.8\Us$，波的频率为 $1.25\UHz$，B 正确；从图像中可知 $\lambda=4\Um$，根据公式得 $v=\lambda/T=5\Ums$，故 A 错误；$x$ 坐标为 $15\Um$ 的质点和 $x$ 坐标为 $3\Um$ 的质点相隔 $12\Um$，为波长的整数倍，即两质点为同相点，而 $x$ 坐标为 $3\Um$ 的质点经过 $t=0.2\Us$ 即四分之一周期振动到平衡位置，所以 $x$ 坐标为 $15\Um$ 的质点在 $t=0.2\Us$ 时振动到平衡位置，C 错误；$x$ 坐标为 $22\Um$ 的质点和 $x$ 坐标为 $2\Um$ 的质点为同相点，$x$ 坐标为 $2\Um$ 的质点经过 $t=0.2\Us$ 即四分之一周期恰好位于波峰，故 $x$ 坐标为 $22\Um$ 的质点在 $t=0.2\Us$ 时恰好位于波峰，D 正确；当质点 P 位于波峰时，经过了半个周期，而 $x$ 坐标为 $17\Um$ 的质点和 $x$ 坐标为 $1\Um$ 的质点为同相点，经过半个周期 $x$ 坐标为 $1\Um$ 的质点恰好位于波谷，E 正确；选项 BDE 正确。

	（2）在半圆柱形玻璃砖横截面内，考虑沿半径方向射到圆心 $O$ 的光线 1（如图），它在圆心处的入射角 $\theta_1$ 为 $\theta_1=\gamma$①，恰好等于临界角，发生全反射。

	在光线 1 左侧的光线（例如光线 2），经柱面折射后，射在玻璃砖底面上的入射角 $\theta_2$ 满足 $\theta_2>\gamma$②，因而在底面上发生全反射，不能直接折射出。

	在光线 1 右侧的光线（例如光线 3），经柱面折射后，射在玻璃砖底面上的入射角 $\theta_3$ 满足 $\theta_3<\gamma$③，因而在底面上不能发生全反射，能从玻璃砖底面射出。

	射到半圆柱面最右侧的光线 4 与柱面相切，入射角 $i$ 为 $i=\frac{\pi}{2}$④，

	由折射定律知，经圆柱面折射后的折射角 $\angle OAB=\theta_4$ 满足 $\sin i=n\sin\theta_4$⑤，

	式中，$n$ 是玻璃的折射率。由全反射角的定义知 $1=n\sin\gamma$⑥，

	联立④⑤⑥式得 $\theta_4=\gamma$⑦，

	由几何关系知 $\angle AOB=\gamma$，故底面上透光部分的宽度 $OB$ 为 $l=\frac{R}{2\cos\gamma}$。

	\onepicture[h=4.0cm]{figs/16c.jpg}
}

\item
%% number: 17
%% paperName: 2015年海南理综第 17 题 8 分
%% typeId: 6
%% type: 计算题
%% chapter: 运动和力
%% point: 动量和能量
%% method: 无
%% score: 8
%% degree: 450
%% duplicateId: 0
%% body:
【物理——选修 3-5】（12 分）

（1）（4 分）氢原子基态的能量为 $E_1=-13.6\UeV$。大量氢原子处于某一激发态。由这些氢原子可能发出的所有光子中，频率最大的光子能量为 $-0.96E_1$，频率最小的光子的能量为\tkanswer[1.4cm]{0.31}$\UeV$（保留 2 位有效数字），这些光子可具有\tkanswer[1.0cm]{10} 种不同的频率。

（2）（8 分）运动的原子核 $\ce{^{A}_{Z}X}$ 放出 $\alpha$ 粒子后变成静止的原子核 Y。已知 X、Y 和 $\alpha$ 粒子的质量分别是 $M$、$m_1$ 和 $m_2$，真空中的光速为 $c$，$\alpha$ 粒子的速度远小于光速。求反应后与反应前的总动能之差以及 $\alpha$ 粒子的动能。

%% answer:
\jdanswer{
\begin{enumerate}
	\item
	$0.31\UeV$，10 种

	\item
	$(M-m_1-m_2)c^2$，$\frac{M}{M-m_2}(M-m_1-m_2)c^2$
\end{enumerate}
}
%% memo:
\memoanswer{
	（1）频率最大的光子能量为 $-0.96E_1$，即 $E_n-(-13.6\UeV)=-0.96\times(-13.6\UeV)$，得 $E_n=-0.54\UeV$。又 $E_n=E_1/n^2$，故 $n=5$；从 $n=5$ 能级开始，共有 $5\to1$、$5\to2$、$5\to3$、$5\to4$、$4\to1$、$4\to2$、$4\to3$、$3\to1$、$3\to2$、$2\to1$，10 种不同频率的光子。频率最小的光子的能量应为 $5\to4$ 跃迁时发出的光子，且 $E_5=-0.54\UeV$、$E_4=-0.85\UeV$，所以 $\Delta E=E_5-E_4=0.31\UeV$。

	（2）反应后由于存在质量亏损，所以反应前后总动能之差等于质量亏损而释放出的能量，故根据爱因斯坦质能方程可得 $\frac{1}{2}m_2v_\alpha^2-\frac{1}{2}Mv_X^2=(M-m_1-m_2)c^2$；反应过程中三个粒子组成的系统动量守恒，故有 $Mv_X=m_2v_\alpha$，联立可得 $\frac{1}{2}m_2v_\alpha^2=\frac{M}{M-m_2}(M-m_1-m_2)c^2$。
}

\end{enumerate}

\end{document}
